\section{Scope, method and the layer model}
\label{sec:scope}

\subsection{Scope}

We include a technology if it (i) is a published specification, standard or
widely adopted open vocabulary; (ii) carries \emph{machine-interpretable meaning}
beyond byte-level syntax; and (iii) is used, or is plausibly usable, in
discrete-manufacturing, process, energy or building automation. We therefore
include information-modelling standards that are not usually called ``semantic''
(\opcua, AutomationML, PackML) because they encode exactly the kind of meaning an
agent needs, and we exclude pure transport protocols (Modbus, PROFINET RT,
EtherCAT) except where they are cited as the semantic baseline against which
everything else is an improvement.

We exclude: proprietary vendor data models without public specification;
machine-learning feature stores, which are statistical rather than semantic
artefacts; and general enterprise data catalogues, except where they carry a
formal vocabulary (DCAT).

\subsection{Method}

Specifications were collected from the standardisation bodies themselves (IEC,
ISO, ISA, VDI/VDE/NAMUR, W3C, ETSI, OPC~Foundation, IDTA, ECLASS~e.V.,
AutomationML~e.V., IDSA, Eclipse Foundation), from the reference documentation
of the corresponding open implementations, and from the peer-reviewed literature
introducing the major domain ontologies. The ten AUF criteria
(\S\ref{sec:framework}) were then derived top-down from a failure taxonomy of
industrial agents (Table~\ref{tab:failure-taxonomy}): each criterion exists
because at least one observed failure mode is caused by its absence. Each
technology was scored $0$--$3$ per criterion against the rubric in
Appendix~\ref{app:rubric}, on the basis of what the \emph{specification mandates
or standardises}, not what a particular implementation happens to offer.

Finally, the framework's central claim, that the criteria are independent and
non-substitutable, is testable, and we test it by ablation
(\S\ref{sec:design}). If the criteria were substitutable, adding more of one
(say, type information) would compensate for the absence of another (say, unit
semantics). Experiments E2, E4 and E5 show that it does not.

\paragraph{A note on scoring.} AUF scores are an expert assessment, not a
measurement, and we mark them as such. The experimental sections make no use of
them; they are a navigational aid for the reader and a summary of the prose.
Where a technology is extensible in principle but silent by default (\opcua{}
alarm limits, AAS qualifiers) we score the default, since that is what an agent
receives in practice.

\subsection{The layer model}
\label{sec:layers}

Discussions of ``semantics'' in automation founder because the word is used for
five different things. We separate them explicitly; the separation is the
backbone of the survey and of the experiment.

\begin{table*}[t]
\centering
\caption{The five-layer model used throughout this survey. Each layer answers a
different question, and no layer subsumes the one below it.}
\label{tab:layers}
\footnotesize
\setlength{\tabcolsep}{4pt}
\begin{tabular}{@{}L{1.9cm}L{2.6cm}L{6.0cm}L{4.4cm}@{}}
\toprule
\textbf{Layer} & \textbf{Question answered} & \textbf{Representative technologies} &
\textbf{Agent affordance} \\
\midrule
L0 Identification &
\emph{Which concept is this?} &
ECLASS / IEC~CDD (IEC~61360) / IEC~61987, IRDI (ISO~29002-5),
GS1~GTIN/GIAI &
Stable, global, resolvable reference; cross-vendor property equality \\[3pt]

L1 Information models &
\emph{What is it, and what can be done to it?} &
\opcua{}~/~IEC~62541 (+ companions), AAS (IDTA), AutomationML~/~IEC~62714,
MTP (VDI/VDE/NAMUR 2658), PackML~/~ISA-TR88.00.02, ISA-95~/~IEC~62264,
W3C~WoT~TD, NGSI-LD, IEC~61850 &
Types, cardinalities, ranges, access rights, callable methods, state machines \\[3pt]

L2 Formal semantics &
\emph{What follows, and what is forbidden?} &
RDF, RDFS, OWL~2, SHACL, SPARQL, PROV-O &
Entailment, constraint validation, path queries, aggregation, provenance \\[3pt]

L3 Domain ontologies &
\emph{What does this mean in \emph{this} domain?} &
QUDT, SOSA/SSN, Brick, SAREF (+4BLDG, 4INMA, 4ENER), DEXPI &
Quantity kinds and unit calculus, observation semantics, equipment taxonomies,
process semantics \\[3pt]

L4 Dataspaces &
\emph{May I use it, and can I believe it?} &
IDS-RAM \& Dataspace Protocol, Catena-X~/~Manufacturing-X, ODRL &
Usage policy, sovereignty, attestation, lineage across organisational boundaries \\
\bottomrule
\end{tabular}
\end{table*}

Three properties of this model matter for the argument that follows.

\paragraph{Layers are complements, not alternatives.} A common
mis-specification (``we will use the AAS instead of an ontology'', or ``we
have \opcua, we do not need RDF'') treats L1 and L2/L3 as competing choices.
They answer different questions. \opcua{} tells an agent that
\texttt{PT0301.PV} is a \texttt{Double} with an \texttt{EURange} of
$[0, 6]$ and an \texttt{EUInformation} display name \texttt{"bar"}. It does not
tell the agent that \texttt{bar} and \texttt{psi} are commensurable and how, that
the variable is downstream of a pasteuriser, or that writing its sibling
setpoint is currently forbidden. Those are L3, L2 and L1-behavioural facts
respectively, and the experiment in \S\ref{sec:results} shows each of them
independently moving the numbers.

\paragraph{Adoption is inverted with respect to agent value.} L0 and L1 are the
mature, widely deployed layers; L2 and L3 are patchy; L4 is nascent. The AUF
matrix (Table~\ref{tab:auf-matrix}) shows that the agent-critical criteria of
unit calculus, behaviour, constraints and provenance concentrate in exactly the
layers with the thinnest deployment.

\paragraph{The layers have very different token economics.} L1 exchange formats
are verbose by design, because they were built for machines with disks, not for
models with context windows. Experiment~E3 (Table~\ref{tab:e3-serialisation})
measures the same plant at $1\times$ as a historian CSV and up to
$\MaxSerialRatio\times$ as expanded RDF, which forces a design consequence
explored in \S\ref{sec:integration}: the semantic layer must be
\emph{queried}, never pasted.

\subsection{Terminology}

We use \emph{tag} or \emph{signal} for a process variable (an \opcua{}
\texttt{VariableNode}, an AAS \texttt{Property}, a historian point);
\emph{asset} for a physical thing with an identity; \emph{grounding} for the
mapping of a natural-language expression onto an identified entity;
\emph{context} for the token sequence placed in the model's window;
\emph{guardrail} for any mechanism that accepts or rejects a proposed agent
action before execution; and \emph{ceiling} for the maximum task accuracy
achievable by a reader that never invents facts (\S\ref{sec:design}).
